\documentclass[conference]{IEEEtran}
\usepackage{cite}
\usepackage{amsmath,amssymb}
\usepackage{booktabs}
\usepackage{graphicx}
\usepackage{array}
\usepackage{url}
\usepackage{hyperref}

\title{MudraVision: Real-Time Bharatanatyam Mudra Recognition Using Transfer Learning With MobileNetV2 and Deployment-Oriented Computer Vision Integration}

\author{\IEEEauthorblockN{Rishitha et al.}
\IEEEauthorblockA{Mudra Recognition Project\\
India\\
Email: contact@example.com}}

\begin{document}
\maketitle

\begin{abstract}
Hand mudra recognition is an important problem for culturally grounded human-computer interaction, assistive interfaces, and performance analytics in Indian classical dance. This paper presents an end-to-end mudra recognition system trained on the Bharatanatyam Mudra Dataset and deployed for real-time operation through OpenCV, MediaPipe, Flask, and a web frontend. The proposed learning pipeline uses grayscale input at 224x224 resolution, transfer learning with an ImageNet-pretrained MobileNetV2 backbone, two-stage training (frozen-head training followed by selective fine-tuning), and augmentation tailored for webcam variability. The notebook-derived dataset contains 28,431 images across 50 mudra classes, split with a training/validation partition of 80/20. The final reported validation accuracy is 96.34\%. A full-class evaluation pass reports 97.29\% overall accuracy with macro-averaged F1-score of 0.9718 and weighted F1-score of 0.9728. Beyond model training, this work emphasizes deployment realism: live hand localization, optional HSV/YCrCb segmentation modes, smoothed prediction stabilization, and method-selectable runtime launching from a frontend.
\end{abstract}

\begin{IEEEkeywords}
Mudra recognition, Bharatanatyam, transfer learning, MobileNetV2, hand gesture recognition, real-time inference, MediaPipe, computer vision deployment.
\end{IEEEkeywords}

\section{Introduction}
Hand gestures are among the most information-rich nonverbal modalities in human communication. In Bharatanatyam and related Indian classical traditions, mudras encode symbolic, narrative, and expressive semantics. Automatic recognition of such gestures is relevant for dance pedagogy, archival indexing, performer feedback, culturally contextual HCI, and multimodal assistive systems.

Despite strong progress in generic hand gesture recognition, mudra recognition presents specific challenges: fine-grained inter-class similarity, high intra-class variation across performers and viewpoints, illumination variability in webcam conditions, and deployment constraints for low-latency user-facing inference. To address these, we design and evaluate a transfer-learning system centered on MobileNetV2, using grayscale training and augmentation to improve robustness while preserving deployment efficiency.

The core contributions are:
\begin{enumerate}
\item A reproducible notebook-to-deployment pipeline for 50-class mudra recognition.
\item A two-stage optimization strategy combining frozen-backbone learning and controlled fine-tuning.
\item A practical deployment stack integrating OpenCV, MediaPipe, Flask, and frontend controls.
\item Comprehensive class-wise metrics and deployment-focused analysis.
\end{enumerate}

\section{Related Work}
Vision-based gesture recognition has progressed from handcrafted descriptors to deep learning systems \cite{lecun2015deep}. Lightweight architectures such as MobileNet are especially relevant for real-time deployment \cite{sandler2018mobilenetv2,howard2019mobilenetv3}. Transfer learning improves data efficiency and stability \cite{pan2010transfer,weiss2016survey}. Landmark-based real-time tracking frameworks such as MediaPipe are widely adopted for practical systems \cite{bazarevsky2020blazepose}. Data augmentation remains essential for robustness under viewpoint and lighting changes \cite{perez2017augmentation,shorten2019survey}.

In sign and gesture recognition literature, region extraction, deep classification, and temporal stabilization are common design patterns \cite{koller2016deep,simonyan2014twostream}. Our work follows this direction while focusing on Bharatanatyam mudras and deployment integration.

\section{Dataset and Problem Formulation}
\subsection{Dataset Source}
The training notebook clones the Bharatanatyam-Mudra-Dataset repository \cite{mudradataset}. After cleaning hidden/non-class artifacts and normalizing folder names, the pipeline identifies:
\begin{itemize}
\item Total classes: 50
\item Total images: 28,431
\end{itemize}

\subsection{Data Split}
The notebook uses:
\begin{itemize}
\item Validation split: 0.2
\item Seed: 123
\item Grayscale loading with \texttt{image\_dataset\_from\_directory}
\end{itemize}
Resulting partitions:
\begin{itemize}
\item Training: 22,745 images
\item Validation: 5,686 images
\end{itemize}

\subsection{Task Definition}
Given an image \(x\), classify it into one of \(K=50\) mudra classes using single-label multi-class classification.

\section{Methodology}
\subsection{Input Pipeline}
Training uses grayscale images resized to \(224 \times 224\). Since MobileNetV2 expects 3 channels, grayscale input is replicated:
\begin{equation}
X_{3c} = \mathrm{concat}(X_g, X_g, X_g)
\end{equation}
followed by standard MobileNetV2 preprocessing.

\subsection{Augmentation}
On-the-fly augmentation includes horizontal flip, small rotation, zoom, translation, and contrast perturbation.

\subsection{Model Architecture}
The architecture is:
\begin{enumerate}
\item Input \(224 \times 224 \times 1\)
\item Data augmentation
\item Channel replication to 3 channels
\item MobileNetV2 backbone (\texttt{include\_top=False}, ImageNet weights)
\item GlobalAveragePooling2D
\item BatchNormalization
\item Dropout (0.35)
\item Dense softmax layer (50 classes)
\end{enumerate}

Total parameters from notebook summary: 2,327,154.

\subsection{Two-Stage Training}
\textbf{Stage 1 (frozen backbone):}
\begin{itemize}
\item Backbone non-trainable
\item Adam, learning rate \(1\times10^{-3}\)
\item Up to 12 epochs
\end{itemize}

\textbf{Stage 2 (fine-tuning):}
\begin{itemize}
\item Backbone unfrozen after layer index 100
\item Adam, learning rate \(1\times10^{-5}\)
\item Early stopping and checkpointing
\end{itemize}

Loss function:
\begin{equation}
\mathcal{L} = -\frac{1}{N}\sum_{i=1}^{N}\log p(y_i|x_i)
\end{equation}

\section{Deployment Architecture}
\subsection{System Components}
Deployment consists of:
\begin{itemize}
\item \texttt{original.py}: core recognizer and hand cropper
\item \texttt{skin\_segmenation.py}: HSV/YCrCb segmentation and live prediction UI
\item \texttt{app.py}: Flask integration layer with method-trigger endpoints
\item React frontend: user-facing model selection controls
\end{itemize}

\subsection{Runtime Flow}
\begin{figure}[t]
\centering
\fbox{\parbox{0.95\columnwidth}{
\textbf{Runtime Pipeline Diagram}\\
Frontend Button (YCrCb/HSV) $\rightarrow$ Flask \texttt{/api/run-segmentation}\\
$\rightarrow$ Launch \texttt{skin\_segmenation.py} with selected method\\
$\rightarrow$ OpenCV Webcam Stream\\
$\rightarrow$ MediaPipe Hand Detection\\
$\rightarrow$ Segmentation (HSV or YCrCb)\\
$\rightarrow$ MobileNetV2 Inference\\
$\rightarrow$ Top-1 / Top-3 / Hasta Type display
}}
\caption{Deployment flow used by the integrated system.}
\label{fig:runtime}
\end{figure}

\subsection{Method-Specific Launch Behavior}
The frontend now calls a backend endpoint that launches:
\begin{itemize}
\item YCrCb: \texttt{python skin\_segmenation.py --camera 0 --method ycrcb}
\item HSV: \texttt{python skin\_segmenation.py --camera 0 --method hsv}
\end{itemize}

\section{Experimental Results}
\subsection{Validation Metrics}
Notebook evaluation reports:
\begin{itemize}
\item Validation loss: 0.1197
\item Validation accuracy: 0.9634 (96.34\%)
\end{itemize}

\subsection{Full Evaluation Pass}
A full directory evaluation block reports:
\begin{itemize}
\item Overall accuracy: 0.9729
\item Macro precision: 0.9725
\item Macro recall: 0.9718
\item Macro F1-score: 0.9718
\item Weighted F1-score: 0.9728
\end{itemize}

\subsection{Class-Wise Behavior}
Several classes achieved 100\% recall (e.g., Berunda, Chakra, Garuda, Katakavardhana, Matsya, Shivalinga). More difficult classes include Bramaram (84.62\%), Ardhapathaka (85.10\%), and Tripathaka (86.03\%), indicating confusion among visually similar mudras.

\begin{table}[t]
\caption{Quantitative Summary}
\label{tab:summary}
\centering
\begin{tabular}{>{\raggedright\arraybackslash}p{0.58\columnwidth}p{0.30\columnwidth}}
\toprule
Metric & Value \\
\midrule
Number of classes & 50 \\
Total images & 28,431 \\
Training split & 22,745 \\
Validation split & 5,686 \\
Validation accuracy & 96.34\% \\
Full-eval accuracy & 97.29\% \\
Macro F1-score & 0.9718 \\
Weighted F1-score & 0.9728 \\
Model parameters & 2,327,154 \\
\bottomrule
\end{tabular}
\end{table}

\section{Discussion}
Performance indicates that lightweight transfer learning can successfully model fine-grained mudra classes under realistic capture conditions. Two-stage training appears to stabilize adaptation from generic ImageNet features to mudra-specific discriminative patterns. Integration with classical segmentation enables runtime flexibility in varying backgrounds and illumination conditions.

An important caveat is that one notebook evaluation step computes predictions on the full image directory, which includes images seen during train/validation construction. Therefore, that value should be interpreted as a broad diagnostic metric rather than a strict unseen-test benchmark. Future protocol design should include a disjoint held-out test set and, ideally, subject-disjoint splits.

\section{Limitations and Future Work}
\subsection{Current Limitations}
\begin{enumerate}
\item No independently curated test split is reported.
\item Cross-subject generalization is not explicitly measured.
\item No confidence calibration or uncertainty quantification is reported.
\item Dynamic gesture sequence modeling is not included.
\end{enumerate}

\subsection{Future Directions}
\begin{enumerate}
\item Subject-disjoint evaluation with confidence intervals.
\item Comparative study with EfficientNet-Lite and compact ViT variants.
\item Landmark + appearance fusion for ambiguous class pairs.
\item Distillation and quantization for edge deployment.
\item Sequence-aware extensions for expressive dance phrase recognition.
\end{enumerate}

\section{Conclusion}
This paper presented a complete mudra recognition workflow from notebook training to real-time deployment. Using MobileNetV2 transfer learning on a 50-class Bharatanatyam mudra dataset, the system achieved strong reported performance and practical runtime behavior. The integration of model training, segmentation options, live webcam inference, and frontend-triggered execution provides a usable and extensible foundation for culturally grounded gesture interfaces.

\begin{thebibliography}{20}
\bibitem{lecun2015deep}
Y. LeCun, Y. Bengio, and G. Hinton, ``Deep learning,'' \emph{Nature}, vol. 521, no. 7553, pp. 436--444, 2015.

\bibitem{sandler2018mobilenetv2}
M. Sandler, A. Howard, M. Zhu, A. Zhmoginov, and L.-C. Chen, ``MobileNetV2: Inverted residuals and linear bottlenecks,'' in \emph{Proc. CVPR}, 2018, pp. 4510--4520.

\bibitem{howard2019mobilenetv3}
A. Howard \emph{et al.}, ``Searching for MobileNetV3,'' in \emph{Proc. ICCV}, 2019, pp. 1314--1324.

\bibitem{pan2010transfer}
S. J. Pan and Q. Yang, ``A survey on transfer learning,'' \emph{IEEE Trans. Knowl. Data Eng.}, vol. 22, no. 10, pp. 1345--1359, 2010.

\bibitem{weiss2016survey}
K. Weiss, T. M. Khoshgoftaar, and D. Wang, ``A survey of transfer learning,'' \emph{J. Big Data}, vol. 3, no. 9, 2016.

\bibitem{bazarevsky2020blazepose}
V. Bazarevsky \emph{et al.}, ``BlazePose: On-device real-time body pose tracking,'' arXiv:2006.10204, 2020.

\bibitem{gonzalez2018dip}
R. Gonzalez and R. Woods, \emph{Digital Image Processing}, 4th ed. Pearson, 2018.

\bibitem{yang1998face}
M.-H. Yang and N. Ahuja, ``Detecting human faces in color images,'' in \emph{Proc. ICIP}, 1998, pp. 127--130.

\bibitem{perez2017augmentation}
L. Perez and J. Wang, ``The effectiveness of data augmentation in image classification using deep learning,'' arXiv:1712.04621, 2017.

\bibitem{shorten2019survey}
A. Shorten and T. M. Khoshgoftaar, ``A survey on image data augmentation for deep learning,'' \emph{J. Big Data}, vol. 6, no. 60, 2019.

\bibitem{koller2016deep}
O. Koller, H. Ney, and R. Bowden, ``Deep hand: How to train a CNN on 1 million hand images when your data is continuous and weakly labelled,'' in \emph{Proc. CVPR}, 2016, pp. 3793--3802.

\bibitem{simonyan2014twostream}
K. Simonyan and A. Zisserman, ``Two-stream convolutional networks for action recognition in videos,'' in \emph{Proc. NIPS}, 2014, pp. 568--576.

\bibitem{chollet2021dl}
F. Chollet, \emph{Deep Learning with Python}, 2nd ed. Manning, 2021.

\bibitem{everingham2010voc}
M. Everingham \emph{et al.}, ``The Pascal Visual Object Classes (VOC) challenge,'' \emph{Int. J. Comput. Vis.}, vol. 88, no. 2, pp. 303--338, 2010.

\bibitem{mudradataset}
J. Raj R., ``Bharatanatyam-Mudra-Dataset,'' GitHub repository. [Online]. Available: \url{https://github.com/jisharajr/Bharatanatyam-Mudra-Dataset}

\bibitem{tensorflow}
TensorFlow, ``TensorFlow: An end-to-end open-source machine learning platform.'' [Online]. Available: \url{https://www.tensorflow.org/}

\bibitem{opencv}
OpenCV, ``Open Source Computer Vision Library.'' [Online]. Available: \url{https://opencv.org/}

\bibitem{kerasmobilenet}
Keras Team, ``Keras applications: MobileNetV2.'' [Online]. Available: \url{https://keras.io/api/applications/mobilenet/}

\bibitem{mediapipe}
MediaPipe Team, ``MediaPipe solutions and tasks.'' [Online]. Available: \url{https://developers.google.com/mediapipe}

\bibitem{sklearnmetrics}
Scikit-learn Developers, ``Classification metrics and confusion matrix.'' [Online]. Available: \url{https://scikit-learn.org/stable/modules/model_evaluation.html}
\end{thebibliography}

\end{document}
